% ===========================================================================
% Figure: Sensor architecture — pluggable SDK interface
% Shows ISensorSource as an abstract interface with two concrete implementations.
% ===========================================================================
\begin{figure*}[t]
\centering
\begin{tikzpicture}[node distance=1.2cm]

% ---- Abstract interface (top) ----
\node[block, minimum width=6cm, minimum height=1.5cm, fill=blue!15]
  (iface) {\textbf{\texttt{ISensorSource}}\\[2pt]
           \texttt{+ open()} \quad \texttt{+ hasNext()}\\
           \texttt{+ nextFrame()} \quad \texttt{+ close()}};

% ---- Concrete implementations (bottom) ----
\node[blockgreen, below left=1.5cm and 0.5cm of iface, minimum width=4.5cm,
      minimum height=2.2cm]
  (reader) {\textbf{DatasetReader}\\[3pt]
            \footnotesize Reads from disk:\\
            \footnotesize • RGB images (JPG/PNG)\\
            \footnotesize • Depth maps (PNG/EXR)\\
            \footnotesize • LiDAR PLY files\\
            \footnotesize • Poses (TUM/KITTI/EuRoC)\\
            \footnotesize • IMU (CSV)};

\node[blockorange, below right=1.5cm and 0.5cm of iface, minimum width=4.5cm,
      minimum height=2.2cm]
  (orbbec) {\textbf{OrbbecSource}\\[3pt]
            \footnotesize Live camera stream:\\
            \footnotesize • OrbbecSDK v2 (C API)\\
            \footnotesize • Femto Mega / Gemini\\
            \footnotesize • LiDAR + iToF + stereo\\
            \footnotesize • IMU (on-board)\\
            \footnotesize \emph{(compile-gated)}};

% ---- Future extension ----
\node[block, below=3.5cm of iface, minimum width=4.5cm, fill=gray!10, dashed]
  (future) {\textbf{FutureSource} (extensible)\\[2pt]
            \footnotesize Any SDK implementing\\
            \footnotesize \texttt{ISensorSource}};

% ---- Inheritance arrows ----
\draw[dasharrow] (reader.north) -- (iface.south -| reader.north)
  node[midway, right, font=\footnotesize] {\texttt{implements}};
\draw[dasharrow] (orbbec.north) -- (iface.south -| orbbec.north)
  node[midway, left, font=\footnotesize] {\texttt{implements}};
\draw[dasharrow] (future.north) -- (iface.south)
  node[midway, right, font=\footnotesize] {\texttt{implements}};

% ---- Group ----
\begin{scope}[on background layer]
  \node[group=Sensor Source Layer, fit=(iface)(reader)(orbbec)(future),
        inner sep=14pt, fill=blue!2] {};
\end{scope}

\end{tikzpicture}
\caption{The \texttt{ISensorSource} abstract interface and its concrete
implementations. The \texttt{DatasetReader} handles offline file-based data,
while the \texttt{OrbbecSource} wraps the Orbbec SDK for live camera input.
The interface is designed for extensibility: any future sensor SDK can be
integrated by implementing four virtual methods.}
\label{fig:sensor_architecture}
\end{figure*}
